\section{A coprime carrier: thirteen integers at degree 1088}
\label{sec:crt1088}

The preceding construction can be reduced further without adding a
variable or a fixed coefficient. The explicit factor \(q^3\) in the
scale \(Y=(2v+1)q^3\) is expensive after repeated substitution. We instead
supply \(Y=qv\), then recover the missing divisibility and odd quotient
through the same root norm that recovers the main Pell congruence.
The two relevant moduli are provably coprime on every integer assignment.
This gives an explicit polynomial of exact degree \(1088\) in the same
thirteen unrestricted integer unknowns.

\subsection{Complete specification}

The coordinates remain
\[
j,C,z,h,r,v,w,T,c,n,t,p_2,p_3\in\Z.
\]
Use the same six fixed numerals \(\beta,K_0,M_C,M_F,
A_{\rm cone},B_{\rm cone}\) as in Section~\ref{sec:equation}; in
particular \(\beta=2^{1064921}-1\), and \(M_F\) is already shifted.
The following definitions replace the earlier front end:
\begin{align}
J&=6j+1,&q&=\beta J+1,&W&=12w+4,\nonumber\\
X&=qW,&Y&=qv,&E&=XY,\nonumber\\
a&=E+Y,&H&=4a+3,&\Delta&=a^2+4a+3.
\label{eq:crtfront}
\end{align}
The transport parametrization is unchanged:
\begin{align}
F&=(K_0+W)C-z(q-1),&\alpha&=q-F-C,\nonumber\\
R&=\bigl(q(q-F)-C\bigr)(q^2-1)+(M_C+qM_F)J,\nonumber\\
k&=R+1+hE,&\eta&=c-Yk,&\zeta&=k-\eta.
\label{eq:crttransport}
\end{align}
Put
\begin{align}
i&=2r+1,&I&=ic^2,&\Df&=1+\Delta I^2,\nonumber\\
Q&=\Delta^2I^2,&\delta&=Q-1,&V&=\Df(cT-R)-c,\nonumber\\
C_y&=QV^2-1,&\Dt&=1+XY^2(XY^2+1)k^2,\nonumber\\
\Dm&=1+\Delta c^2,&b_0&=X+ac.
\label{eq:crtpell}
\end{align}
The rational notation \(\Dy=C_y/\delta\) is used only in the proof.
Define
\begin{align}
G&=(A_{\rm cone}F)(B_{\rm cone}C)
   -(q-A_{\rm cone}F-B_{\rm cone}C)^2,\nonumber\\
\Ag&=v^2\bigl[G\eta\zeta-(t^2-G-\eta\zeta+1)^2\bigr]
       -1-p_2^2-p_3(p_3+1).
\label{eq:crtguard}
\end{align}
The multiplier \(v^2\) in the guard is intentional. It will force
\(v\ne0\), and its effect on the radical's magnitude will be controlled.

For the new carrier set
\begin{align}
M&=2q^2,&\Lambda&=4H+1,&U&=MH,\nonumber\\
Z&=U(n+V)+M\Lambda b_0+H(v-q^2),\nonumber\\
A_0&=M^2\Lambda^2\Dm,&A_1&=U^2\Df,&A_2&=U^2\Dt,\nonumber\\
A_3&=U^2\Ag,&A_4^{\rm num}&=U^2C_y.
\label{eq:crtcarrier}
\end{align}
Here \(M\) is a carrier modulus, unrelated to the native mask numerals.
The requested equation is the single ordinary polynomial
\begin{equation}
\boxed{\quad
P_{13}^{\rm CRT}
=\delta^{16}\F_{32}\left(
Z;A_0,A_1,A_2,A_3,\frac{A_4^{\rm num}}{\delta}
\right)=0.
\quad}
\label{eq:crtpolynomial}
\end{equation}
The denominator is cleared symbolically, exactly as in
Section~\ref{sec:circuit}. No additional existential variable, variable
exponent, numerical root extraction, or division operation is part of
the polynomial.

\subsection{Unconditional realness, including the zero-scale branch}

At every integer assignment, \(q\) is even,
\(q\equiv2\pmod3\), and \(|q|\ge16\). The factor \(3w+1\)
never vanishes, so \(|X|\ge64\). The new \(Y=qv\) may initially
vanish; it must not be silently assumed nonzero.

If \(v=0\), then \(Y=a=0\), \(\Delta=3\), and \(H=3\).
If \(v\ne0\), then \(|Y|\ge16\) and
\(|a|=|X+1||Y|\ge63\cdot16\), so \(\Delta>0\) and
\(|H|>3\). In both cases \(\Delta\) is odd, \(H\ne0\), and
\begin{equation}
3<\frac{\Lambda}{H}=4+\frac1H<5.
\label{eq:crtlambda}
\end{equation}
In particular \(\Lambda\ne0\), and the carrier scales are nonzero.

The radicands \(\Dm\) and \(\Df\) are at least one. With
\(u=XY^2\in\Z\), the integer inequality \(u(u+1)\ge0\) gives
\(\Dt=1+k^2u(u+1)\ge1\). For the auxiliary radicand, if \(c=0\),
then \(Q=0,\delta=-1,C_y=-1\), and \(\Dy=1\). If \(c\ne0\),
then
\[
Q=\Delta^2i^2c^4>1,\qquad \Df=1+\Delta i^2c^4>|c|+1.
\]
Writing \(b=cT-R\in\Z\), either \(b=0\) and \(|V|=|c|\ge1\),
or \(b\ne0\) and \(|V|\ge\Df-|c|>1\). Therefore
\[
\Dy=V^2+\frac{V^2-1}{Q-1}\ge1.
\]
We have proved, without using a guard, a Pell equation, or native
decoding, that \(\delta\ne0\) and all four non-guard radicands are
positive on every integer assignment.

At a zero of \eqref{eq:crtpolynomial}, one of the conjugate linear
factors vanishes over \(\mathbb C\). If \(\Ag<0\), its radical
would be the only nonzero imaginary summand in a factor with a real
center and four real summands. Thus \(\Ag\ge0\). It follows that
\begin{equation}
v^2\Pair(G,\eta\zeta,t)
=\Ag+1+p_2^2+p_3(p_3+1)>0.
\label{eq:crtguardrecover}
\end{equation}
Consequently \(v\ne0\) and \(\Pair(G,\eta\zeta,t)>0\). The
joint-positivity lemma gives \(G>0\) and \(\eta\zeta>0\).
The latter excludes \(c=0\): that case would give
\(\eta\zeta=-Y(Y+1)k^2\le0\). Hence at a zero we have
\(c\ne0,Q>1,\delta>0\).

This proves the required hypotheses in their proper order. In
particular, \(v=0\) is excluded because its guard radicand is strictly
negative; it was not excluded when proving unconditional realness.

\subsection{Separation with only \texorpdfstring{$|Y|\ge|q|$}{|Y|>=|q|}}

Set
\[
x=|X|,\qquad m_Y=|Y|,\qquad K=|k|,\qquad C_0=|c|.
\]
The recovered strict guards imply
\begin{equation}
x\ge64,\quad m_Y\ge|q|\ge16,\quad K\ge2,\qquad
K(m_Y-1)<C_0<K(m_Y+1).
\label{eq:crtweakbounds}
\end{equation}
Indeed \(\eta/k\in(0,1)\) and \(c=k(Y+\eta/k)\).
In particular \(C_0>30\). Let
\[
d_+=\sqrt{\Dm},\quad f_+=\sqrt{\Df},\quad
\tau_+=\sqrt{\Dt},\quad y_+=\sqrt{\Dy},\quad
g_+=\sqrt{\Ag},\quad \lambda=\Lambda/H.
\]
The guard root \(g_+\) may be zero; the others are positive.

As before, the determinant inequality \(G>0\) proves \(G<q^2\):
with \(U_c=A_{\rm cone}F\), \(W_c=B_{\rm cone}C\), its two factors
have the same sign and \(|U_c|+|W_c|<2|q|\).
Since \(\eta\zeta\le k^2/4\), the scaled guard gives
\begin{equation}
g_+<\frac{|v||q|K}{2}
=\frac{Km_Y}{2}<\frac{8C_0}{15}.
\label{eq:crtguardbound}
\end{equation}
The last bound uses \(m_Y\ge16\) and
\(C_0>K(m_Y-1)\).

The first norm gives
\begin{equation}
K(xm_Y^2-1)<\tau_+<K(xm_Y^2+1),\qquad \tau_+>2C_0.
\label{eq:crttaubound}
\end{equation}
The main root has a useful lower comparison with \(\tau_+\):
\[
d_+>\bigl((x-1)m_Y-3\bigr)K(m_Y-1)>\frac34\tau_+.
\]
For the second inequality, expansion gives
\begin{align*}
&4\bigl((x-1)m_Y-3\bigr)(m_Y-1)-3(xm_Y^2+1)\\
&\hspace{12mm}=m_Y^2(x-4)-4m_Y(x+2)+9>0
\end{align*}
for \(x\ge64,m_Y\ge16\). Combining this with
\eqref{eq:crtlambda} deliberately places the scaled main root above
the first root:
\begin{equation}
\lambda d_+>2\tau_+.
\label{eq:crtmainbound}
\end{equation}
Also
\[
f_+^2>\Delta C_0^4,\qquad
d_+^2=\Delta C_0^2+1<2\Delta C_0^2,
\]
so
\begin{equation}
f_+>\frac{C_0}{2}d_+>15d_+>2\lambda d_+.
\label{eq:crtstrongbound}
\end{equation}

The estimates \(f_+>3C_0^2\) and \(Q>f_+^2\) hold here too,
since \(|a|\) is large and \(i\ne0\). Thus the same two auxiliary
cases remain valid:
\begin{equation}
\begin{cases}
C_0<y_+<C_0+1,&cT-R=0,\\
y_+>2f_+,&cT-R\ne0.
\end{cases}
\label{eq:crtycases}
\end{equation}
In the first case,
\(y_++g_+<(23/15)C_0+1<2C_0<\tau_+\).
Equations~\eqref{eq:crtmainbound} and \eqref{eq:crtstrongbound}
give strict superincreasing order
\[
f_+,\quad \lambda d_+,\quad \tau_+,\quad y_+,\quad g_+.
\]
In the second case the order is
\[
y_+,\quad f_+,\quad \lambda d_+,\quad \tau_+,\quad g_+.
\]
Zero \(g_+\) is omitted. Each entry exceeds the sum of all entries
to its right, not merely the next entry.

Dividing a vanishing factor by \(U\) expresses a rational number
as a signed sum of precisely these magnitudes. Lemma~\ref{lem:radicals}
makes all five roots rational. The roots \(d,f,\tau,p_1\) of the
integer radicands \(\Dm,\Df,\Dt,\Ag\) are integral, and
Lemma~\ref{lem:descent} makes the rational auxiliary root \(y\)
integral. The factor's signs can be absorbed into these integer roots.

\subsection{Two coprime congruences, with no hidden scale promise}

Because \(Y=qv\),
\[
H=4(X+1)qv+3\equiv3\pmod q.
\]
The fixed odd exponent gives \(\beta\equiv1\pmod3\), and
\(J=6j+1\), so \(q\equiv2\pmod3\). Since \(H\) is odd,
\begin{equation}
\gcd(H,M)=\gcd(H,2q^2)=1,\qquad \Lambda\equiv1\pmod H.
\label{eq:crtcoprime}
\end{equation}
These are universal arithmetic facts about the displayed formulas.
They are not assertions about unspecified canonical witnesses.

The signed integer roots satisfy
\begin{equation}
U(n+V)+M\Lambda b_0+H(v-q^2)
=M\Lambda d+U(f+\tau+p_1+y).
\label{eq:crtrootrelation}
\end{equation}
Modulo \(M\), this gives \(H(v-q^2)\equiv0\), so
\eqref{eq:crtcoprime} yields
\begin{equation}
v=q^2(2v_{\rm old}+1),\qquad
v_{\rm old}=\frac{v-q^2}{2q^2}\in\Z.
\label{eq:crtoldscale}
\end{equation}
Consequently \(Y=qv=(2v_{\rm old}+1)q^3\), exactly the original
source scale, including its odd quotient. Modulo \(H\),
\eqref{eq:crtrootrelation} gives
\(M\Lambda(d-b_0)\equiv0\). Both multipliers are units modulo
\(H\), so
\begin{equation}
\gamma=\frac{d-b_0}{H}\in\Z.
\label{eq:crtoldgamma}
\end{equation}
This restores the main Pell root and its required congruence.

The normalized strong and auxiliary witnesses are restored by
\eqref{eq:restore}--\eqref{eq:restore2}, unchanged. Since
\(\Pair(G,\eta\zeta,t)>0\), the universal ternary \(\pos\)
representation supplies the old private guard variables. They need
not equal the new ones: the latter represent
\(v^2\Pair(G,\eta\zeta,t)\). This distinction is necessary for
soundness. The transport equation already holds identically by
\eqref{eq:crttransport}.

\begin{theorem}[Coprime-carrier reduction]\label{thm:crt1088}
The polynomial \(P_{13}^{\rm CRT}\) has an integer zero if and only
if the normalized integer system \(\Cal\) has an integer solution.
Consequently, with the same inherited ordinary calibration,
\[
\mathrm{ZFC}\ \vdash\
\left(\neg\exists\boldsymbol{x}\in\Z^{13}
\ P_{13}^{\rm CRT}(\boldsymbol{x})=0\right)
\ \longleftrightarrow\ \Con(\mathrm{ZFC}).
\]
The polynomial has exact total degree \(1088\).
\end{theorem}
\begin{proof}
The preceding argument restores a solution of \(\Cal\) from a
polynomial zero. Conversely, start with a solution of \(\Cal\),
whose old scale is \(Y=(2v_{\rm old}+1)q^3\), and put
\[
v=Y/q=q^2(2v_{\rm old}+1)\ne0.
\]
All native quantities remain unchanged. Retain its \(t\), for which
\(\Pair(G,\eta\zeta,t)>0\), and regenerate \(p_1,p_2,p_3\)
as a \(\pos\) representation of the positive integer
\(v^2\Pair(G,\eta\zeta,t)\). Then the new \(\Ag=p_1^2\).
Finally choose
\begin{equation}
n=\Lambda\gamma+f+\tau+p_1+y-V-\frac{v-q^2}{M}\in\Z.
\label{eq:crtchoose-n}
\end{equation}
Every quantity on the right is fixed before \(n\), and substitution
gives \eqref{eq:crtrootrelation}. One norm factor therefore vanishes.
This is the exact sign convention for the displayed centered carrier.
The calibration conclusion composes this new algebraic equivalence
with the inherited ordinary proof for \(\Cal\). Exact degree is
proved below.
\end{proof}

For a canonical refutation witness the same construction uses its
actual \(Y\); no extra divisibility or freedom of choice is asserted.
The earlier odd strong normalization is also unchanged: with
\(A=a+2\) even and \(R,c\) odd, the recurrence gives
\(\psi_A(Rc)\) odd, and division by the odd \(c^2\) gives odd
\(i\). This parity fact follows from the recurrence, not from only
the first term of a binomial congruence.

\subsection{Exact degree and a coefficient certificate}

The lower-degree scale changes the relevant degrees to
\[
\begin{array}{c|rrrrrrrr}
\text{register}&a&\Delta&\Df&Q&R&V&C_y&\Ag\\\hline
\text{degree}&4&8&14&22&5&19&60&30
\end{array}
\]
We also have \(\deg\Dt=22\), \(\deg U=6\), and
\(\deg Z=25\). In the biquadratic evaluator of
Section~\ref{sec:circuit}, take
\[
d=M^2\Lambda^2\Dm,\quad e=U^2\Ag,\quad
A=U^2\Df,\quad B=U^2\Dt,\quad \Cn=U^2C_y.
\]
Their respective degrees are \(22,42,26,34,72\). Assign formal
weights \(11\) and \(21\) to the roots of \(d\) and \(e\).
Writing
\[
s=Un+M\Lambda b_0+H(v-q^2),\qquad Z=s+UV,
\]
the scalar coefficient of the first norm register has the exact cut
\begin{equation}
J_0=\delta(s^2+2UVs+d+e-A-B)-U^2(V^2-1).
\label{eq:crtcut}
\end{equation}
Here \(\deg s=11\). The four weighted component degrees are
\[
64,\quad58,\quad68,\quad54.
\]
The unique highest component is \(2\delta Z\sqrt e\), of
weight \(68\). The scalar subtraction in the next squaring has
degree at most \(128<136\), and the last tail subtraction has
weight at most \(248<272\). Therefore the cleared three-root
norm has unique leading scalar
\[
16Q_0^4Z_0^4e_0^2
\]
of degree \(272\). The two quadratic norm operations give exact
degree \(1088\) and top form \(2^{16}Q_0^{16}Z_0^{16}e_0^8\).

For a numerical-slice certificate, define
\begin{align*}
q_0&=6\beta j,&W_0&=12w,&Y_0&=q_0v,&X_0&=q_0W_0,\\
a_0&=X_0Y_0,&i_0&=2r,&F_0&=W_0C-q_0z,&R_0&=-q_0^3F_0,\\
k_0&=R_0+ha_0,&\eta_0&=-Y_0k_0,&Q_0&=a_0^4i_0^2c^4,\\
V_0&=-a_0^2i_0^2c^4R_0,&U_0&=8q_0^2a_0,&Z_0&=U_0V_0,\\
e_0&=-U_0^2v^2\eta_0^4.&&&&
\end{align*}
Direct substitution gives the nonzero leading polynomial
\begin{equation}
\bigl(P_{13}^{\rm CRT}\bigr)_{1088}
=2^{112}q_0^{64}a_0^{128}i_0^{64}c^{128}
 R_0^{16}v^{16}\eta_0^{32}.
\label{eq:crttop}
\end{equation}
In particular,
\begin{equation}
\left[j^{480}w^{160}v^{208}h^{32}r^{64}c^{128}z^{16}\right]
P_{13}^{\rm CRT}
=2^{976}3^{640}\beta^{480}>0.
\label{eq:crtcoefficient}
\end{equation}
The exponent \(h^{32}\) uniquely selects the \(ha_0\) term in
\(k_0\) throughout \(\eta_0^{32}\); \(z^{16}\) then uniquely
selects the \(q_0^4z\) term throughout \(R_0^{16}\).
All remaining factors are monomials. The exponent sum is \(1088\),
so this proves exact degree on the actual fixed coefficient slice,
independently of possible relations among the other five numerals.

The gain comes from a complete carrier construction with a proved
coprimality identity and explicit two-way witness maps. The inherited calibration and its stated formalization boundary remain
exactly those identified in the earlier sections.
